\section{公司法体系 - 知识结构和发展逻辑}
\begin{enumerate}
  \item 55 分。发起人、公司、股东他们之间的行为、关系\bs{不能割裂开来单独看}，不然某些题会做错。
  \item 商主体法。
        \begin{enumerate}
          \item 公司法。半壁江山。
                \begin{enumerate}
                  \item 股东的有限责任、公司的独立法人地位。
                  \item 人和钱。股东\zd{出资制度} \yjt 形式、要求、责任。\zd{股东资格} \yjt 名册、登记、冒名、名义。
                  \item 成立后 \yjt 公司分类、\zd{能力}（投资、担保）、\zd{变更}（注册资本、人员、合并、分立）。
                  \item \zd{股东权利}。分红、知情权、优先认购、直接诉讼、代表诉讼。
                  \item 3 个板块。人与机构（董高监）、股权和股份（流转）、特殊情况（上市公司）。
                \end{enumerate}
          \item 合伙企业。
          \item 生命周期。设立、运行、变更、消亡。
        \end{enumerate}
  \item 商行为法。
        \begin{enumerate}
          \item 规范退出。破产 \yjt 清算、\zd{重整}（股权、债务、业务调整）、和解。
          \item 权益性工具。证券法 \yjt 股票 + 债券 + 基金。
          \item 避险工具。保险法 \yjt 人身 + 财产。
          \item 便利支付。票据法 \yjt 支票 + 汇票 + 本票。
          \item 信托法。
        \end{enumerate}
\end{enumerate}
